% Part 3: representing the past (working title).
% Build from this directory: make pt3
% In progress: motivation, HMM intuition, cave, and formal definition.
% Target a lecture comparable in length to Part 2. Students have completed
% the Part 2 generator exercises. A reading session on Shai et al. follows.
\documentclass[aspectratio=169,12pt,t]{beamer}
\usepackage{iliad-slides}
\usepackage{tikz}
\usetikzlibrary{arrows.meta,shapes.geometric,fit}
\graphicspath{{figures/}}
% Factor colours carried over from the original HMM slides.
\definecolor{hmmInitial}{RGB}{253,246,124}
\definecolor{hmmFirst}{RGB}{150,235,150}
\definecolor{hmmSecond}{RGB}{142,233,246}
\definecolor{hmmLast}{RGB}{250,206,152}
\definecolor{hmmSum}{RGB}{248,132,132}
\newcommand{\hmmfactor}[2]{\colorbox{#1}{\ensuremath{#2}}}
\title{Computational Mechanics}
\subtitle{3. Representing the past}
\institute{ILIAD Intensive}
\date{Autumn 2026}
\author{Xavier Poncini (Simplex)}

\begin{document}
\iliadtitle

\begin{frame}{Outline}
  \vspace{5mm}
  \begin{enumerate}
    \setlength{\itemsep}{4mm}
    \item Hidden Markov models
    \item Belief states and next-token prediction
    \item Sufficient statistics and belief-state minimality
    \item Generalised HMMs and minimal-dimensional realisations
  \end{enumerate}
\end{frame}

\begin{frame}{Representing the past for prediction}
  \setlength{\parskip}{0pt}
  \vfill
  Last lecture, we established that RNNs and transformers that achieve
  \textbf{exact next-token prediction} develop sufficient statistics for the
  \textbf{entire future}.

  \vspace{9mm}
  In this lecture, we will study \textbf{hidden Markov models} and their
  generalisations. We will construct belief states and next-token vectors,
  then determine when belief states are \textbf{minimal sufficient statistics}.

  \vspace{9mm}
  The answer will suggest the kinds of internal representations that neural
  networks trained on such data might learn.
  \vfill

  \note{The recap retains the recurrent-state assumptions from Part 2.
  Sufficiency concerns the model's full predictive state, rather than an
  arbitrary activation at a single layer or position. Exact prediction is
  required across supported histories. The lecture constructs belief states,
  distinguishes them from next-token vectors, and then uses minimal GHMM
  realisations to study when belief states are minimal sufficient statistics.
  This result motivates, but does not prove, the hypothesis that neural
  networks trained on these processes learn analogous internal
  representations. The reading session examines evidence for that hypothesis.}
\end{frame}

\begin{frame}{What is a hidden Markov model?}
  \setlength{\parskip}{0pt}
  \vfill
  Intuitively, a \textbf{hidden Markov model} (HMM) is a stochastic process
  that consists of two components:

  \vspace{5mm}
  \begin{itemize}
    \setlength{\itemsep}{5mm}
    \item A \textbf{hidden}, unobserved component that evolves according to
      a \textbf{Markovian dynamic}.
    \item An \textbf{observed} component that provides a possibly lossy view
      of the hidden dynamic.
  \end{itemize}
  \vfill

  \note{This restores the intuitive two-component introduction from slide 3
  of the old Part 2 deck. The cave illustration will follow, then the formal
  Mealy-type definition, word probabilities, and examples.
  The later formal definition specifies a joint distribution over the next
  emitted symbol and destination state conditional on the current state.
  The two components described here need not be conditionally independent.
  Hidden means not directly observed; observations may nevertheless allow
  the current hidden state to be inferred exactly.}
\end{frame}

\begin{frame}
  \vspace*{6mm}
  \vfill
  {\centering\evidenceimage[57mm]{hmm-cave.png}\par}
  \vfill

  \note{Original labelled Plato's cave illustration from the old Part 2
  deck, slide 4. The image is copied without modification, retaining its
  Hidden Dynamic and Observations labels. Use the illustration to connect
  the hidden dynamics to what an observer can see: the observations may
  discard distinctions in the underlying state. This is an analogy for
  partial observation, rather than a specification of the Markov property.
  The formal Mealy-type definition follows on the next slide.}
\end{frame}

\begin{frame}
  \vspace*{6mm}
  \setlength{\parskip}{0pt}
  \setlength{\abovedisplayskip}{4pt}
  \setlength{\belowdisplayskip}{4pt}
  A \textbf{Mealy-type} (edge-emitting) HMM consists of:

  \vspace{2mm}
  \begin{itemize}
    \setlength{\itemsep}{1mm}
    \item A finite set $S$ of \textbf{hidden states}.
    \item A finite alphabet $\mathcal X$ of \textbf{observed symbols}.
    \item A collection of $|S|\times |S|$ \textbf{transition matrices}
      $(T^{(a)})_{a\in\mathcal X}$.
    \item An \textbf{initial distribution} $\eta^{(\emptyset)}$ over $S$.
  \end{itemize}

  \vspace{2mm}
  The matrices specify the joint emission and state transition probabilities:
  \[
    T^{(a)}_{ij}
      :=\Pr\!\left(X=a,\ S'=j\mid S=i\right).
  \]
  \par\vspace{3mm}
  We can conveniently represent these transitions with labelled diagrams:

  \vfill
  {\centering
    \begin{tikzpicture}[
      every node/.style={text=iliadInk,font=\normalsize},
      state/.style={circle,draw=iliadAccent,line width=0.9pt,
        fill=iliadAccent!7,minimum size=9mm,inner sep=1pt},
      flow/.style={-{Stealth[length=2.2mm]},draw=iliadAccent,line width=0.9pt}
    ]
      \node[state] (i) at (0,0) {$i$};
      \node[state] (j) at (3,0) {$j$};
      \draw[flow] (i) -- node[above] {$a:p$} (j);
      \node[anchor=west] at (4,0) {$p=T^{(a)}_{ij}$};
    \end{tikzpicture}\par}
  \vfill

  \source{Moore-type HMMs emit from states; Mealy-type HMMs emit on
  transitions. They describe the same processes, but may require different
  numbers of hidden states.}

  \note{Adapted from the formal definition on slide 5 of the original
  Part 2 deck. State labels use i,j. In the displayed conditional,
  S and S' denote the current and next hidden-state random variables,
  and X denotes the symbol emitted on that transition. The unprimed S
  also names the state set in the opening list; context distinguishes
  the set from the random variable. In time-indexed notation, S_t is
  the state after X_{1:t}, before the transition emitting X_{t+1}.
  Rows index source states and columns index destination states. The
  time-homogeneous joint transition law is independent of the earlier
  history conditional on S_t. Emission and destination need not be
  conditionally independent of each other.

  Each symbol matrix is nonnegative and substochastic; their sum is a
  row-stochastic matrix. The normalization condition sums over both symbols
  and destinations. The initial distribution is a stochastic row vector
  and need not be stationary. It is included here to specify the model
  completely; the following slide develops its role in word probabilities.
  An edge labelled a:p emits a and moves to the indicated state with
  joint probability p.

  Footnote source: Travers and Crutchfield, Equivalence of History and
  Generator Epsilon-Machines, Section II,
  \url{https://csc.ucdavis.edu/~cmg/papers/hgem.pdf}.}
\end{frame}

\begin{frame}
  \vspace*{6mm}
  \setlength{\parskip}{0pt}
  \setlength{\abovedisplayskip}{3pt}
  \setlength{\belowdisplayskip}{3pt}
  \setlength{\fboxsep}{1pt}
  Fix an initial distribution over hidden states, written as a row vector:
  \[
    \eta^{(\emptyset)}_i
      :=\Pr(S_0=i),\qquad i\in S.
  \]

  \vfill
  The probability of observing $x_{1:n}\in\mathcal X^n$ is given by:
  \[
    \begin{aligned}
      \Pr(X_{1:n}=x_{1:n})
        &=\sum\nolimits_{i_0,\ldots,\colorbox{hmmSum}{$\scriptstyle i_n$}\in S}
          \hmmfactor{hmmInitial}{\Pr(S_0=i_0)}\\[1mm]
        &\quad\times\hmmfactor{hmmFirst}{\Pr(X_1=x_1,S_1=i_1\mid S_0=i_0)}\\[1mm]
        &\quad\times\hmmfactor{hmmSecond}{\Pr(X_2=x_2,S_2=i_2\mid S_1=i_1)}\cdots\\[1mm]
        &\quad\times\hmmfactor{hmmLast}{\Pr(X_n=x_n,S_n=i_n\mid S_{n-1}=i_{n-1})}.
    \end{aligned}
  \]

  \vfill
  Matrix multiplication performs these sums over hidden states:
  \[
    \Pr(X_{1:n}=x_{1:n})
      =\hmmfactor{hmmInitial}{\eta^{(\emptyset)}}\,
       \hmmfactor{hmmFirst}{T^{(x_1)}}\,
       \hmmfactor{hmmSecond}{T^{(x_2)}}\cdots
       \hmmfactor{hmmLast}{T^{(x_n)}}\,
       \hmmfactor{hmmSum}{\mathbf1^\top}.
  \]
  \par\vspace{2mm}
  Here $\mathbf1$ is the $|S|$-dimensional row vector of ones.

  \vfill
  \textbf{Key point:} probabilities are generated by a linear dynamic!
  \vfill

  \note{Adapted from slide 6 of the old Part 2 deck, using the sequence
  notation and matching factor highlights from the original HMM slides:
  yellow initial distribution, green first transition, cyan second
  transition, orange final transition, and red terminal summation vector.
  We use the sequence
  notation from the current Part 2. Pr(a,j|i) is shorthand for Pr(X=a,S'=j|S=i)
  from the preceding slide. S_0 denotes the initial state before any
  observations. Hidden states are labelled directly by i, not s_i.

  Conditional on each current hidden state, the next emission and
  destination are independent of earlier history. This gives the product
  for each hidden path; marginalizing over all paths gives the word law.
  Matrix multiplication sums over intermediate states, and the final
  column vector $\mathbf1^\top$ sums over the destination. The j-th component of the
  unnormalized row vector is Pr(X_{1:n}=x_{1:n},S_n=j), a joint probability,
  not yet the posterior conditioned on x_{1:n}. This observation prepares the
  later belief-state definition. No stationarity assumption is required.}
\end{frame}

\begin{frame}
  \vspace*{6mm}
  \setlength{\parskip}{0pt}
  \textbf{Example 1: Biased coin}

  \vspace{2mm}
  Each observation is an independent coin toss, with
  $\Pr(X_t=1)=p$, where $0<p<1$.
  \[
    S=\{0\},\qquad\mathcal X=\{0,1\},\qquad
    \eta^{(\emptyset)}=\begin{bmatrix}1\end{bmatrix}.
  \]
  \vfill
  {\centering
  \begin{tikzpicture}[
    every node/.style={text=iliadInk,font=\normalsize},
    state/.style={circle,draw=iliadAccent,line width=0.9pt,
      fill=iliadAccent!7,minimum size=10mm,inner sep=1pt},
    flow/.style={-{Stealth[length=2.2mm]},draw=iliadAccent,line width=0.9pt}
  ]
    \node[state] (coin) at (0,0) {$0$};
    \draw[flow] (coin.north west)
      .. controls +(-1.8,1.1) and +(-1.8,-1.1) ..
      node[left] {$0:1-p$} (coin.south west);
    \draw[flow] (coin.south east)
      .. controls +(1.8,-1.1) and +(1.8,1.1) ..
      node[right] {$1:p$} (coin.north east);
  \end{tikzpicture}\par}
  \vfill
  {\small
    \textbf{Example output:} $\mathtt{1011001010\ldots}$\par}
  \[
    T^{(0)}=\begin{bmatrix}1-p\end{bmatrix},
    \qquad
    T^{(1)}=\begin{bmatrix}p\end{bmatrix}.
  \]
  \vfill
  \note{First example from slide 7 of the old HMM deck. The single
  hidden state is labelled 0, following the direct state-label convention.
  State labels and emission symbols are separate objects despite sharing
  a glyph. These are one-by-one matrices. Their product gives the usual
  Bernoulli sequence probability. The displayed output is one illustrative
  finite sample. This revisits the IID example from Part 2.}
\end{frame}

\begin{frame}
  \vspace*{6mm}
  \setlength{\parskip}{0pt}
  \setlength{\abovedisplayskip}{4pt}
  \setlength{\belowdisplayskip}{4pt}
  \textbf{Example 2: Zero--One--Random}

  \vspace{2mm}
  The generator cycles through three phases: emit $0$, emit $1$, then
  emit a random symbol with $\Pr(1)=p$, where $0<p<1$.
  \[
    S=\{0,1,R\},\qquad \mathcal X=\{0,1\},\qquad
    \eta^{(\emptyset)}=\begin{bmatrix}\tfrac13&\tfrac13&\tfrac13\end{bmatrix}.
  \]

  \vfill
  {\centering% Edge-emitting Zero--One--Random generator with Bernoulli(p) random phase.
% Editable diagram; labels have the form emitted symbol : edge probability.
\begin{tikzpicture}[
  every node/.style={text=iliadInk,font=\normalsize},
  state/.style={circle,draw=iliadAccent,line width=0.9pt,
    fill=iliadAccent!7,minimum size=9mm,inner sep=1pt},
  flow/.style={-{Stealth[length=2.2mm]},draw=iliadAccent,line width=0.9pt}
]
  \node[state] (zero) at (0,0) {$0$};
  \node[state] (one) at (3.3,0) {$1$};
  \node[state] (random) at (6.6,0) {$R$};
  \draw[flow] (zero) -- node[above] {$0:1$} (one);
  \draw[flow] (one) -- node[above] {$1:1$} (random);
  \draw[flow] (random.north)
    .. controls +(0,0.8) and +(0,0.8) ..
    node[above] {$0:1-p,\quad 1:p$} (zero.north);
\end{tikzpicture}
\par}
  \vfill
  {\small
    \textbf{Example output:} $\mathtt{011\,010\,011\,011\ldots}$
    (starting in phase $0$).\quad State order: $(0,1,R)$.\par}
  \[
    T^{(0)}=\begin{bmatrix}
      0&1&0\\
      0&0&0\\
      1-p&0&0
    \end{bmatrix},
    \qquad
    T^{(1)}=\begin{bmatrix}
      0&0&0\\
      0&0&1\\
      p&0&0
    \end{bmatrix}.
  \]
  \vfill
  \note{Second example from slide 7 of the old HMM deck, using the
  uniform initial phase distribution from the Part 2 exercises. It is
  stationary because the phase advances deterministically around a cycle.
  The two labels on the return edge denote distinct emissions with the
  same destination. Rows are source states, columns are destination states.
  The row sums across both matrices total one for each state.
  The observer sees emitted symbols, not phase boundaries. This is the
  same generator used in the exercise, with s_0,s_1,s_R renamed 0,1,R.
  The displayed output begins in phase 0. Its three-symbol blocks correspond
  to the phases 0,1,R, with the random phase emitting 1,0,1,1 in succession.
  Later slides will revisit this model to calculate beliefs and recover
  the causal-state automaton that students have already constructed.}
\end{frame}

\begin{frame}{Predicting HMM observations}
  \setlength{\parskip}{0pt}
  \setlength{\abovedisplayskip}{7pt}
  \setlength{\belowdisplayskip}{1pt}
  \setlength{\fboxsep}{1pt}
  We have observed $X_{1:n}=x_{1:n}$ and want to predict the next symbol:
  \[
    \Pr(X_{n+1}=a\mid X_{1:n}=x_{1:n}).
  \]
  \vfill
  If we knew the current hidden state $S_n=i$, prediction would be
  straightforward:
  \[
    \Pr(X_{n+1}=a\mid S_n=i)=\sum_{j\in S}T^{(a)}_{ij}.
  \]
  \vfill
  With the hidden state unobserved, we average over its possible values:
  \[
    \Pr(X_{n+1}=a\mid X_{1:n}=x_{1:n})
      =\sum_{i\in S}
        \underbrace{\hmmfactor{hmmInitial}{\Pr(S_n=i\mid X_{1:n}=x_{1:n})}}
          _{\text{belief state component}}
        \sum_{j\in S}T^{(a)}_{ij}.
  \]
  \vfill

  The \textbf{belief state} is the row vector

  \vspace{3mm}
  {\centering
    \setlength{\fboxsep}{2mm}
    \setlength{\fboxrule}{0.6pt}
    \fcolorbox{iliadAccent}{iliadAccent!6}{%
      $\displaystyle
        \eta^{(x_{1:n})}_i
          :=\Pr(S_n=i\mid X_{1:n}=x_{1:n}),\qquad i\in S
      $}
    \par}
  \vfill

  \note{Assume that the observed history has positive probability.
  For a current state with positive posterior mass, the edge-emitting
  Markov property gives
  Pr(X_{n+1}=a|S_n=i,X_{1:n}=x_{1:n})=sum_j T^(a)_{ij}.
  The law of total probability, conditioning throughout on the observed
  history, then gives the final expression. Zero-posterior states contribute
  zero. For unreachable states the transition matrix still specifies the
  generative emission kernel.

  Highlight the posterior as the quantity that changes with the observed
  history. The generator's transition probabilities remain fixed.
  This distribution over the current hidden state is the belief state,
  defined at the bottom of this slide. The following slide derives how to
  compute it. The highlighted quantity is a posterior
  distribution, not an entropy or other scalar measure of uncertainty.}
\end{frame}

\begin{frame}{Computing and updating the belief state}
  \setlength{\parskip}{0pt}
  \setlength{\abovedisplayskip}{4pt}
  \setlength{\belowdisplayskip}{4pt}
  For a history $x_{1:n}$ with positive probability, Bayes' rule gives
  \[
    \eta^{(x_{1:n})}_i
      =\Pr(S_n=i\mid X_{1:n}=x_{1:n})
      =\frac{\Pr(S_n=i,X_{1:n}=x_{1:n})}
      {\Pr(X_{1:n}=x_{1:n})}.
  \]

  \vfill
  Using the HMM transition matrices,
  \[
    \eta^{(x_{1:n})}
      =\frac{\eta^{(\emptyset)}T^{(x_1)}\cdots T^{(x_n)}}
      {\eta^{(\emptyset)}T^{(x_1)}\cdots T^{(x_n)}\mathbf1^\top}.
  \]

  \vfill
  When the next token $x_{n+1}$ is observed, the belief updates recursively:
  \[
    \eta^{(x_{1:n})}
      \mathrel{\overset{x_{n+1}}
        {\mapstochar\relbar\joinrel\longrightarrow}}
    \eta^{(x_{1:n+1})}
      =
    \frac{\eta^{(x_{1:n})}T^{(x_{n+1})}}
    {\eta^{(x_{1:n})}T^{(x_{n+1})}\mathbf1^\top}.
  \]
  \vfill

  The current belief and the newly observed token determine the next belief.
  The full observation history does not need to be revisited.
  \vfill

  \note{This version combines the direct calculation with its recursive
  form. Bayes' rule expresses each component as a joint probability divided
  by the probability of the observed history. The matrix product gives the
  corresponding row-vector expression directly. When a new token arrives,
  the previous belief replaces the initial distribution and earlier matrix
  product. Multiplying by the token-labelled transition matrix gives
  unnormalised probabilities over S_{n+1}; dividing by their sum produces
  the new belief. The denominator is positive whenever the extended history
  is possible. With our timing convention, eta after x_{1:n} is a
  distribution over S_n, and eta after x_{1:n+1} is a distribution over
  S_{n+1}.}
\end{frame}

\begin{frame}{The next-token probability vector}
  \setlength{\parskip}{0pt}
  \setlength{\abovedisplayskip}{4pt}
  \setlength{\belowdisplayskip}{4pt}
  \setlength{\fboxsep}{1pt}
  For a history $x_{1:n}$, define the row vector
  \[
    p^{(x_{1:n})}
      :=\left(
        \Pr(X_{n+1}=a\mid X_{1:n}=x_{1:n})
      \right)_{a\in\mathcal X}.
  \]

  \vfill
  Returning to the expression on slide 10,
  \[
    \hmmfactor{hmmFirst}{\Pr(X_{n+1}=a\mid X_{1:n}=x_{1:n})}
      =\sum_{i\in S}
        \hmmfactor{hmmInitial}{\Pr(S_n=i\mid X_{1:n}=x_{1:n})}
        \hmmfactor{hmmSecond}{\sum_{j\in S}T^{(a)}_{ij}}.
  \]

  \vfill
  Using the belief state and next-token vector, this becomes
  \[
    \hmmfactor{hmmFirst}{p^{(x_{1:n})}}
      =\hmmfactor{hmmInitial}{\eta^{(x_{1:n})}}
       \hmmfactor{hmmSecond}{A},
    \qquad
    A_{ia}:=\sum_{j\in S}T^{(a)}_{ij}.
  \]

  \vfill
  The HMM fixes $A$, so the map from belief states to next-token vectors is
  linear. Distinct belief states can produce the same next-token vector.
  \vfill

  \note{The next-token probability vector records one probability for every
  token in the alphabet. Return explicitly to the law-of-total-probability
  expression on slide 10, now replacing the posterior term with the belief
  state notation. For a fixed hidden state i, the sum over j is the
  probability of emitting token a next. The matrix A stores these
  probabilities, with one row per hidden state and one column per token.
  Therefore multiplying the belief row vector by A gives the full next-token
  vector. Emphasise that A depends only on the HMM and that this fixed linear
  readout can be many-to-one.}
\end{frame}

\begin{frame}{Sufficient statistics for HMM prediction}
  \setlength{\parskip}{0pt}
  \setlength{\abovedisplayskip}{4pt}
  \setlength{\belowdisplayskip}{4pt}
  Recall that a statistic $R$ is \textbf{sufficient for prediction}
  if histories with the same statistic have the same probability for every
  finite continuation:
  \[
    R(x)=R(x')
      \quad\Longrightarrow\quad
    \Pr(y\mid x)=\Pr(y\mid x')
      \qquad\text{for every }y\in\mathcal X^*.
  \]

  \vfill
  \textbf{The next-token probability vector} $p^{(x_{1:n})}$ gives all one-symbol\\
  continuation probabilities, but not necessarily those for longer continuations.\\
  It need not be sufficient for the \textbf{entire future}.
  {\small\color{iliadMuted}\emph{(See exercises for details.)}}

  \vfill
  \textbf{The belief state} $\eta^{(x_{1:n})}$ preserves the probability of
  every finite continuation of an HMM. It is sufficient for predicting the
  \textbf{entire future}.

  \vfill

  \note{The target of prediction matters. The next-token vector is itself a
  minimal sufficient statistic for the one-step prediction problem, but the
  Part 2 definition asks for the distribution of every finite continuation.
  The next-token vector can merge histories that agree about the next symbol
  but disagree about later symbols. The following slide proves that the belief
  state does retain the full continuation law for an HMM.}
\end{frame}

\begin{frame}{The belief state is a sufficient statistic}
  \setlength{\parskip}{0pt}
  \setlength{\abovedisplayskip}{3pt}
  \setlength{\belowdisplayskip}{3pt}
  Given a possible history $x_{1:n}$ and finite continuation $y_{1:m}$,
  Bayes' rule gives
  \[
    \begin{aligned}
      &\Pr(X_{n+1:n+m}=y_{1:m}\mid X_{1:n}=x_{1:n}) \\[2pt]
      &\quad=
      \frac{\Pr(X_{1:n+m}=x_{1:n}y_{1:m})}
           {\Pr(X_{1:n}=x_{1:n})} \\[4pt]
      &\quad=
      \frac{
        \eta^{(\emptyset)}T^{(x_1)}\cdots T^{(x_n)}
        T^{(y_1)}\cdots T^{(y_m)}\mathbf1^\top
      }{
        \eta^{(\emptyset)}T^{(x_1)}\cdots T^{(x_n)}\mathbf1^\top
      } \\[4pt]
      &\quad=
      \eta^{(x_{1:n})}T^{(y_1)}\cdots T^{(y_m)}\mathbf1^\top.
    \end{aligned}
  \]

  \vfill
  The history enters this expression only through its belief state. Hence,
  \[
    \eta^{(x)}=\eta^{(x')}
      \quad\Longrightarrow\quad
    \Pr(y\mid x)=\Pr(y\mid x')
      \qquad\text{for every }y\in\mathcal X^*.
  \]

  \vfill
  The belief state is sufficient. \textbf{What about minimality?}
  \vfill

  \note{Assume throughout that the histories being conditioned on have
  positive probability. The first equality writes the conditional
  probability as a joint probability divided by the probability of the
  observed history. Substitute the HMM matrix products in the numerator and
  denominator. Their common history-dependent prefix, after normalisation,
  is exactly the belief state. The remaining matrix product depends only on
  the proposed continuation. Therefore histories with the same belief state
  assign the same probability to every finite continuation, which is the
  definition of predictive sufficiency recalled on the previous slide.}
\end{frame}

\begin{frame}{Belief states need not be minimal}
  \setlength{\parskip}{0pt}
  \setlength{\abovedisplayskip}{3pt}
  \setlength{\belowdisplayskip}{3pt}
  Consider the two-state HMM
  \[
    \eta^{(\emptyset)}=\left(\tfrac12,\tfrac12\right),
    \qquad
    T^{(0)}=
    \begin{pmatrix}
      \tfrac12 & 0 \\
      \tfrac12 & 0
    \end{pmatrix},
    \qquad
    T^{(1)}=
    \begin{pmatrix}
      0 & \tfrac12 \\
      0 & \tfrac12
    \end{pmatrix}.
  \]

  \vfill
  After observing $0$ or $1$, respectively, the belief state is
  \[
    \eta^{(0)}=(1,0),
    \qquad
    \eta^{(1)}=(0,1).
  \]

  \vfill
  However, the future is a sequence of independent fair coin flips in either
  case:
  \[
    \Pr(y\mid 0)=\Pr(y\mid 1)=2^{-|y|}
      \qquad\text{for every }y\in\{0,1\}^*.
  \]

  \vfill
  Thus, distinct beliefs can contain exactly the same predictive information.
  More generally, define the \textbf{predictive equivalence relation}
  \[
    \eta\sim\eta'
      \quad\Longleftrightarrow\quad
    \eta T^{(y)}\mathbf1^\top
      =\eta'T^{(y)}\mathbf1^\top
      \qquad\text{for every }y\in\mathcal X^*.
  \]

  \vfill
  Beliefs are \textbf{minimal} precisely when no two reachable beliefs are
  \textbf{predictively equivalent}. Generalised HMMs will let us characterise this condition.
  \vfill

  \note{This HMM generates independent fair coin flips. Its hidden state only
  records the most recently observed token: observing 0 places the model in
  the first state, while observing 1 places it in the second. Since the two
  rows of each symbol-labelled transition matrix agree, both hidden states
  induce the same distribution over every future word. The two posterior
  distributions are therefore distinct but predictively equivalent. This
  shows that sufficiency alone does not imply minimality. The equivalence
  relation displayed on the slide identifies beliefs that assign the same
  probability to every finite continuation. Belief states are minimal exactly
  when every such equivalence class contains only one reachable belief. This
  motivates the move to generalised HMMs on the following slide, where
  predictive redundancy can be described in linear-algebraic terms.

  The redundant fair-coin presentation is Upper's example in Section 3.4,
  equation (3.38), of Theory and Algorithms for Hidden Markov Models and
  Generalized Hidden Markov Models (1997).}
\end{frame}

\begin{frame}{Generalised hidden Markov models}
  \setlength{\parskip}{0pt}
  \setlength{\abovedisplayskip}{3pt}
  \setlength{\belowdisplayskip}{3pt}
  A $d$-dimensional \textbf{generalised hidden Markov model} (GHMM) consists of:

  \vspace{1mm}
  \begin{itemize}
    \setlength{\itemsep}{0.5mm}
    \item A finite alphabet $\mathcal X$ of \textbf{observed symbols}.
    \item A collection of $d\times d$ real \textbf{transition matrices}
      $(T^{(a)})_{a\in\mathcal X}$.
    \item An \textbf{initial row vector}
      $\eta^{(\emptyset)}\in\mathbb R^{1\times d}$.
    \item A \textbf{column vector} $\phi\in\mathbb R^{d\times 1}$.
  \end{itemize}

  \vspace{1mm}
  Writing $T=\sum_{a\in\mathcal X}T^{(a)}$, these objects must satisfy
  \[
    T\phi=\phi,\qquad
    \eta^{(\emptyset)}\phi=1,\qquad
    \eta^{(\emptyset)}T^{(w)}\phi\geq 0
      \quad\text{for every }w\in\mathcal X^*.
  \]

  These conditions ensure that the probability of any observed sequence is
  \[
    \Pr(X_{1:n}=x_{1:n})
      =\eta^{(\emptyset)}T^{(x_1)}\cdots T^{(x_n)}\phi.
  \]

  \vspace{1mm}
  An HMM is the special case $d=|S|$ and $\phi=\mathbf1^\top$, where
  $\eta^{(\emptyset)}$ is a probability distribution, each $T^{(a)}$ is
  nonnegative, and $T$ is row-stochastic.
  \vfill

  \note{The definition retains exactly the matrix-product expression used
  for HMM word probabilities, but its internal coordinates no longer need
  to be probabilities. The fixed-point condition T phi equals phi gives
  consistency under extension of a word: summing the probabilities of wa
  over all next symbols a recovers the probability of w. The condition on
  the initial vector normalises the empty word, and nonnegativity of every
  word probability then gives a valid stochastic process.

  The component list deliberately parallels the HMM definition on slide 6.
  A GHMM replaces the finite hidden-state set by a real vector space of
  dimension d and adds the column vector phi used to evaluate word weights.
  An ordinary HMM is the special case in which the symbol-labelled matrices
  are nonnegative, their sum is row stochastic, the initial vector is a
  probability distribution, and phi is the all-ones column vector. In the
  general case, the presentation states are better regarded as basis vectors
  for a linear representation of the process than as physical hidden states.
  Do not call the resulting normalised forward coordinates belief states,
  since they may have negative components.}
\end{frame}

\begin{frame}{Minimality of a GHMM}
  \setlength{\parskip}{0pt}
  \setlength{\abovedisplayskip}{3pt}
  \setlength{\belowdisplayskip}{3pt}
  Histories and continuations generate two linear spaces:

  \[
    \begin{aligned}
      \mathcal H
        &:=\operatorname{span}\!\left\{
          \eta^{(\emptyset)}T^{(x_1)}\cdots T^{(x_n)}
          \;\middle|\; x_{1:n}\in\mathcal X^n,\ n\geq 0
        \right\},\\
      \mathcal F
        &:=\operatorname{span}\!\left\{
          T^{(y_1)}\cdots T^{(y_m)}\phi
          \;\middle|\; y_{1:m}\in\mathcal X^m,\ m\geq 0
        \right\}.
    \end{aligned}
  \]

  A GHMM is \textbf{minimal-dimensional} if no lower-dimensional GHMM
  realises the same process. For a $d$-dimensional GHMM, this is equivalent to
  \[
    \dim\mathcal H=\dim\mathcal F=d.
  \]

  \vfill
  \textbf{Theorem.} If a stochastic process admits a GHMM realisation,
  then it admits a \textbf{minimal-dimensional GHMM realisation}.

  \source{Upper (1997), Theorem 4.3.10.}

  \note{Upper's Theorem 4.3.10 constructs, from any GHMM, an equivalent
  standard presentation based on minimal history and future wordlists and
  proves that no equivalent GHMM has fewer states. Thus every process that
  admits a finite-dimensional GHMM presentation admits a minimal-dimensional
  one.

  In linear-realisation language, reachability and observability identify the
  two possible sources of redundant dimension. If the history space is
  smaller than the ambient row space, restrict the representation to the
  history space. If a nonzero row direction pairs to zero with every vector
  in the future space, quotient out that invisible direction. Conversely,
  when both spaces have dimension d, the process Hankel matrix has rank d, so
  no lower-dimensional GHMM can realise the same process. This is distinct
  from minimality of a predictive statistic. See Upper, Theory and Algorithms
  for Hidden Markov Models and Generalized Hidden Markov Models, Theorem
  4.3.10 and Sections 4.2--4.3.}
\end{frame}

\begin{frame}{When HMM belief states are minimal}
  \setlength{\parskip}{0pt}
  \setlength{\abovedisplayskip}{3pt}
  \setlength{\belowdisplayskip}{3pt}
  If an HMM is minimal-dimensional among all GHMM realisations of its process,
  then its belief state is a \textbf{minimal sufficient statistic} for
  prediction.

  \vspace{2mm}
  Suppose two supported histories $x_{1:n}$ and $x'_{1:n'}$ predict every
  continuation identically.
  Then, for every $y_{1:m}$,
  \[
    \left(
      \eta^{(x_{1:n})}-\eta^{(x'_{1:n'})}
    \right)
    T^{(y_1)}\cdots T^{(y_m)}\mathbf1^\top
    =0.
  \]

  \vspace{1mm}
  Minimal-dimensionality implies that the future vectors
  \[
    T^{(y_1)}\cdots T^{(y_m)}\mathbf1^\top
  \]
  span the entire column space. Therefore,
  \[
    \eta^{(x_{1:n})}=\eta^{(x'_{1:n'})}.
  \]

  \vfill
  Thus, histories have the same belief state precisely when they have the
  same distribution over the future. The belief state is
  \textbf{minimal sufficient}.

  \vspace{-1mm}
  \source{Based on Upper (1997), Sections 4.2--4.3.}
  \vfill

  \note{Sufficiency gives one implication immediately: equal beliefs imply
  equal probabilities for all future continuations. For the converse, the
  displayed difference of two beliefs annihilates every future vector. Since
  the HMM is minimal-dimensional even within the larger class of GHMMs, the
  preceding slide gives dim F equals d. The future vectors therefore span the
  entire column space, so the difference must be zero. This proves that no two
  distinct reachable beliefs are predictively equivalent.}
\end{frame}

\begin{frame}{Two notions of minimality}
  \setlength{\parskip}{0pt}
  \setlength{\abovedisplayskip}{4pt}
  \setlength{\belowdisplayskip}{4pt}
  \small
  Here we summarise how each notion of minimality enters the key result.

  \vspace{2mm}
  \begin{itemize}
    \setlength{\itemsep}{2.5mm}
    \item \textbf{Minimal sufficient statistics:} coarsest partition of the allowed sequences.
    \item \textbf{Minimal-dimensional:} no equivalent GHMM has lower
      dimension.
  \end{itemize}

  \vspace{3mm}
  {\centering
    $\text{HMM is minimal-dimensional}
    \Longrightarrow
    \text{belief states are minimal sufficient statistics}.$\par}

  \vspace{2.5mm}
  Belief states (or their analogues) for a minimal-dimensional GHMM are the
  \textbf{lowest-dimensional} minimal sufficient statistics among those from
  which the entire future distribution is \textbf{linearly decodable}.

  \vspace{2.5mm}
  \textbf{Working hypothesis:} neural networks trained on GHMM data represent
  these objects. {\color{iliadMuted}\itshape
    (We will examine supporting evidence in the reading session.)}

  \vfill

  \note{Predictive minimality is minimality of a statistic for the future
  prediction problem. A minimal sufficient statistic is equivalent to the
  causal-state partition up to relabelling. Minimal-dimensionality instead
  compares all GHMM realisations of the process and asks for the smallest
  linear state-space dimension. The previous slide proves that the stronger
  linear condition makes the reachable belief map injective with respect to
  future laws. The converse can fail because injectivity is required only on
  the set of reachable beliefs, rather than on every direction of the ambient
  linear state space. The closing sentence uses lowest-dimensional only
  within the class of GHMM, or equivalently finite-dimensional linear,
  realisations. It is not a claim about arbitrary numerical encodings of the
  minimal sufficient statistic.}
\end{frame}

\begin{frame}{Summary}
  \setlength{\parskip}{0pt}
  \vfill
  \begin{itemize}
    \setlength{\itemsep}{4mm}
    \item A \textbf{belief state} is a sufficient statistic for prediction:
      the posterior distribution over an HMM's hidden states.
    \item The \textbf{next-token vector} need not be sufficient. It is a
      linear readout of the belief state,
      $p^{(x_{1:n})}=\eta^{(x_{1:n})}A$.
    \item GHMMs generalise HMMs by relaxing the entrywise probabilistic
      constraints on the transition matrices, while retaining the same
      matrix-product expression for output probabilities.
    \item Beliefs need not be minimal. If the HMM is
      \textbf{minimal-dimensional} as a GHMM, they are \textbf{minimal sufficient
      statistics}.
  \end{itemize}

  \vfill

  \note{The first two points distinguish the full predictive information in
  the belief state from its next-token readout. The third recalls that GHMMs
  keep the HMM matrix-product formula while allowing internal coordinates and
  matrices that need not have an entrywise probabilistic interpretation. The
  final point states the main minimality result: minimal-dimensionality in the
  larger GHMM class makes the HMM belief state a minimal sufficient statistic.}
\end{frame}

% Diagram-based summary variant retained for possible later use.
\iffalse
\begin{frame}{Summary}
  \setlength{\parskip}{0pt}
  \vspace{-1mm}
  \centering
  \resizebox{0.86\textwidth}{!}{%
  \begin{tikzpicture}[
    x=1cm,y=1cm,
    every node/.style={text=iliadInk},
    heading/.style={font=\small\bfseries},
    caption/.style={font=\scriptsize,text=iliadMuted},
    hmmstate/.style={circle,draw=iliadAccent,line width=0.8pt,
      fill=iliadAccent!8,minimum size=6.5mm,inner sep=0.5pt,
      font=\scriptsize},
    belief/.style={ellipse,draw=iliadAccent,line width=0.7pt,
      fill=hmmFirst!55,minimum width=11mm,minimum height=5.5mm,
      inner sep=0.5pt,font=\tiny},
    point/.style={circle,draw=iliadAccent,line width=0.7pt,
      minimum size=5mm,inner sep=0pt,font=\tiny},
    flow/.style={-{Stealth[length=1.8mm]},draw=iliadInk,line width=0.7pt},
    map/.style={-{Stealth[length=2.2mm]},draw=iliadMuted,line width=1pt}
  ]
    % Generative realisation.
    \node[heading] at (1.7,4.75) {HMM realisation};
    \node[hmmstate] (h0) at (1.0,3.80) {$0$};
    \node[hmmstate] (h1) at (2.4,3.80) {$1$};
    \node[hmmstate] (h2) at (1.7,2.90) {$2$};
    \draw[flow] (h0) to[bend left=18] (h1);
    \draw[flow] (h1) to[bend left=18] (h0);
    \draw[flow] (h0) -- (h2);
    \draw[flow] (h2) -- (h1);

    % Reachable belief-state presentation.
    \node[heading] at (10.65,4.75) {Reachable belief states};
    \node[belief] (b0) at (10.65,4.20) {$\eta^{(\emptyset)}$};
    \node[belief] (b1) at (9.30,3.55) {$\eta^{(0)}$};
    \node[belief] (b2) at (12.00,3.55) {$\eta^{(1)}$};
    \node[belief,fill=hmmSecond!55] (b3) at (8.70,2.75) {$\eta^{(00)}$};
    \node[belief,fill=hmmSecond!55] (b4) at (10.15,2.75) {$\eta^{(01)}$};
    \node[belief,fill=hmmSecond!55] (b5) at (11.60,2.75) {$\eta^{(10)}$};
    \node[belief,fill=hmmSecond!55] (b6) at (13.05,2.75) {$\eta^{(11)}$};
    \draw[flow] (b0) -- (b1);
    \draw[flow] (b0) -- (b2);
    \draw[flow] (b1) -- (b3);
    \draw[flow] (b1) -- (b4);
    \draw[flow] (b2) -- (b5);
    \draw[flow] (b2) -- (b6);

    \draw[map] (3.05,3.78) -- node[above,font=\scriptsize]
      {Bayesian filtering} (7.85,3.78);

    % Geometry of the same reachable beliefs.
    \node[heading,anchor=west] at (0.05,2.05) {Belief-state geometry};
    \coordinate (t1) at (1.10,0.20);
    \coordinate (t2) at (4.65,0.20);
    \coordinate (t3) at (2.88,1.75);
    \draw[iliadMuted,line width=0.8pt] (t1)--(t2)--(t3)--cycle;
    \node[point,fill=hmmSecond!55] at (1.10,0.20) {$00$};
    \node[point,fill=hmmSecond!55] at (4.65,0.20) {$11$};
    \node[point,fill=hmmSecond!55] at (2.88,1.75) {$01$};
    \node[point,fill=hmmFirst!55] at (2.05,0.95) {$0$};
    \node[point,fill=hmmFirst!55] at (3.55,0.55) {$1$};
    \node[point,fill=white] at (2.75,0.55) {$\emptyset$};

    \draw[map] (8.45,2.58) -- node[above,sloped,font=\scriptsize]
      {same states, plotted} (4.75,1.65);

    % Linear next-token readout.
    \node[heading] at (10.85,1.88) {Next-token geometry};
    \draw[iliadMuted,line width=0.9pt] (8.45,0.62) -- (13.25,0.62);
    \node[point,fill=hmmSecond!55] at (8.70,0.62) {$00$};
    \node[point,fill=hmmSecond!55] at (9.45,0.62) {$01$};
    \node[ellipse,draw=iliadAccent,line width=0.7pt,fill=hmmFirst!55,
      minimum width=10mm,minimum height=5mm,inner sep=0pt,font=\tiny]
      at (10.45,0.62) {$0,\emptyset$};
    \node[point,fill=hmmFirst!55] at (11.50,0.62) {$1$};
    \node[point,fill=hmmSecond!55] at (12.95,0.62) {$11$};
    \node[caption] at (10.85,0.12) {the readout may merge distinct beliefs};

    \draw[map] (5.00,0.78) -- node[above,font=\scriptsize]
      {$p^{(x_{1:n})}=\eta^{(x_{1:n})}A$} (8.05,0.78);
  \end{tikzpicture}
  }

  {\small
    If an HMM is \textbf{minimal-dimensional as a GHMM}, its beliefs are
    \textbf{minimal sufficient statistics}.}

  {\color{iliadMuted}\itshape\scriptsize
    (Reading session: do neural networks trained on GHMM data represent this geometry?)}

  \note{This is a typeset update of the hand-drawn summary diagram in the
  earlier Part 3 deck. The top row moves from a generative HMM or GHMM
  realisation to the reachable beliefs obtained by Bayesian filtering. The
  lower-left panel plots those same beliefs in their state space. The linear
  map A sends belief states to next-token probability vectors; because the
  map can merge distinct beliefs, next-token probabilities need not retain
  all information about the future. The final line adds the minimality result
  established in the present version of the lecture: when the HMM is
  minimal-dimensional as a GHMM realisation, its reachable beliefs are
  minimal sufficient statistics. The reading question connects this
  geometry to representations learned from GHMM data.}
\end{frame}
\fi

% Content plan: add slides as each topic is developed.

% Hidden Markov models
  % Follow the old Part 2 order: intuitive definition (above), Plato's cave,
  % formal Mealy-type definition with edge notation, initial distribution
  % and word probabilities, then biased-coin and Zero--One--Random examples.
  % Reuse figures/z1r-generator.tex when developing the later example slide.
  % Introduce edge-emitting generators, an initial distribution, and T^(a).
  % Preserve the exercises' timing: S_t is the state after observing X_{1:t}.
  % Use Zero--One--Random as the main worked belief-state example later.

% Belief states
  % Define the posterior over S_{|w|} for a supported finite history w.
  % Derive its matrix expression, recursive update, and future-word readout.
  % Prove predictive sufficiency. Calculate Zero--One--Random beliefs and
  % connect its mixed-state presentation to the previous exercise automaton.
  % Plot beliefs in the simplex and distinguish them from next-token outputs.
  % Revisit the Simple Nonunifilar Source as a second worked example.

% Minimality and representation dimension
  % Characterize minimal sufficiency by distinguishability of reachable
  % beliefs through future-word probabilities. Give a redundant counterexample.
  % Distinguish predictive equivalence, HMM state count, and affine dimension.
  % State the representation class whenever making dimensional claims.
  % Retain the finite-history convention, including initial/transient beliefs.

% Belief geometry in transformer activations
  % Prepare the reading session on Shai et al. (NeurIPS 2024):
  % https://arxiv.org/abs/2405.15943
  % Explain the hypothesis, HMM training data, exact belief targets, and
  % fitted linear/affine readout. Compare theoretical and recovered geometry.
  % Explain what a plotted point represents and the role of multiple layers.
  % Distinguish empirical linear decodability from mathematical sufficiency.

% Optional supplementary material, excluded from this initial scaffold:
% GHMM predictive coordinates, normalized updates, and minimal linear
% realizations unique up to similarity. Protect time for the paper discussion.

\end{document}
